\documentclass[12pt,a4paper]{article}
\usepackage{fontspec}
\setmainfont{Times New Roman}
\setsansfont{Arial}
\setmonofont{Consolas}[Scale=0.85]
\usepackage[left=2.6cm,right=2.3cm,top=2.4cm,bottom=2.4cm]{geometry}
\usepackage{amsmath,amssymb,booktabs,array,tabularx}
\usepackage{setspace,enumitem,xcolor,fancyhdr,xurl}
\usepackage[colorlinks=true,linkcolor=blue!45!black,urlcolor=blue!45!black]{hyperref}

\definecolor{navy}{HTML}{0B2239}
\definecolor{teal}{HTML}{157A7F}
\newcommand{\filepath}[1]{\nolinkurl{#1}}
\newcommand{\cue}[1]{\par\smallskip\noindent\textit{\textcolor{teal}{#1}}\par\smallskip}
\newcommand{\slide}[3]{
  \clearpage
  \section*{Slide #1 -- #2}
  \addcontentsline{toc}{section}{Slide #1 -- #2}
  \noindent\textbf{Thời lượng gợi ý:} #3
  \par\medskip
}
\newcommand{\answer}[2]{\paragraph{#1} #2}

\setlength{\parindent}{0.7cm}
\setlength{\emergencystretch}{2em}
\setlength{\headheight}{16pt}
\onehalfspacing
\setlist{nosep,leftmargin=*}
\pagestyle{fancy}
\fancyhf{}
\fancyhead[L]{\small Kịch bản thuyết trình 18--20 phút -- NASA C-MAPSS FD001--FD004}
\fancyhead[R]{\small Phân tích dữ liệu lớn}
\fancyfoot[C]{\thepage}

\begin{document}

\hypersetup{pageanchor=false}
\begin{titlepage}
\centering
{\Large\bfseries KỊCH BẢN THUYẾT TRÌNH\par}
\vspace{0.5cm}
{\large Báo cáo cuối kỳ môn Phân tích dữ liệu lớn\par}
\vspace{1.4cm}
{\LARGE\bfseries Dự đoán tuổi thọ còn lại và cảnh báo sớm\\cho động cơ turbofan\par}
\vspace{0.4cm}
{\Large NASA C-MAPSS FD001--FD004\par}
\vspace{1.5cm}
\begin{tabular}{ll}
Học viên thực hiện: & Phan Nguyễn Mai Phương\\
Mã số học viên: & 25002711\\
Ngành: & Khoa học máy tính
\end{tabular}
\vfill
{\large Tháng 10 năm 2026\par}
\end{titlepage}

\hypersetup{pageanchor=true}

\tableofcontents

\section*{Hướng dẫn sử dụng}
\addcontentsline{toc}{section}{Hướng dẫn sử dụng}
Tài liệu này bám theo deck 20 slide tại
\filepath{slides/cmapss_fd001_presentation.pptx}. Tổng thời lượng phần nói
khoảng 17 phút; demo dashboard khoảng 3 phút, nên toàn bộ trình bày nằm trong
18--20 phút. Speaker notes trong deck chứa bản rút gọn của từng mục dưới đây.

\begin{center}
\begin{tabular}{lll}
\toprule
Phần & Slide & Thời lượng \\
\midrule
Bìa & 1 & 0:20 \\
1. Bài toán và dữ liệu & 2--5 & 2:45 \\
2. Pipeline Big Data & 6--7 & 2:00 \\
3. Phân tích chuỗi thời gian & 8--11 & 2:30 \\
4. Mô hình và đánh giá & 12--16 & 5:00 \\
5. Mở rộng và kết luận & 17--20 & 4:25 \\
\midrule
Tổng phần nói & 20 & 17:00 \\
Demo dashboard & -- & 3:00 \\
\bottomrule
\end{tabular}
\end{center}

Slide làm tròn số liệu đến hai chữ số thập phân; kịch bản dùng số đầy đủ theo
artifact. HGB169 là predictor tham chiếu có điều kiện trong alignment replay;
HGB126 là baseline lịch sử và dashboard legacy. Warning policy được đánh giá
riêng với predictor. Slide 18 dùng window 10 cycles cho FD001--FD004; nhóm metric
này phải được nói tách biệt với panel alignment window 30.

\slide{1}{Bìa}{20 giây}
\cue{Mở deck ở slide 1.}
Kính chào thầy/cô. Em là Phan Nguyễn Mai Phương. Hôm nay em trình bày báo cáo
cuối kỳ môn Phân tích dữ liệu lớn về dự đoán Remaining Useful Life, viết tắt là
RUL, và cảnh báo sớm cho động cơ turbofan trên bộ dữ liệu NASA C-MAPSS.

Mục tiêu là xây dựng một pipeline có thể kiểm tra và tái lập từ raw data đến dự
đoán. Bốn con số ở dưới là khung của bài: 100 train engines, 21 sensor mỗi
cycle, 169 features của HGB169 và MAE trajectory 17,0009 cycles trên panel alignment. Đây là predictor candidate có điều kiện; HGB126 và test metric lịch sử được giữ làm
baseline audit.

\slide{2}{Nội dung và câu hỏi nghiên cứu}{25 giây}
\cue{Chỉ lần lượt năm phần, sau đó chuyển sang khung câu hỏi.}
Bài trình bày có năm phần: bài toán và dữ liệu; pipeline Big Data; phân tích
chuỗi thời gian; mô hình và đánh giá; cuối cùng là phần mở rộng và kết luận.

Sáu câu hỏi nghiên cứu dẫn dắt toàn bài: FD001 có hợp lệ về schema, khóa và
cycle liên tục không; sensor thay đổi thế nào giữa đầu và cuối vòng đời; Pandas
và PySpark có cho cùng kết quả ETL không; HGB126 hay HGB169 ổn định hơn trên cùng
panel; RUL có chuyển thành cảnh báo sớm được không; và model FD001 hoạt động thế
nào khi áp dụng sang FD002--FD004. Ở slide 19, em trả lời lại Q1--Q6 bằng số
liệu.

\slide{3}{Bài toán}{35 giây}
\cue{Giải thích ví dụ bên trái trước, sau đó đọc bảng đối chiếu.}
RUL là số cycle vận hành còn lại trước failure. Đây là bài toán supervised
regression: input là chuỗi 21 sensor theo cycle, output là một số.

Ví dụ trên một train engine: engine kết thúc ở cycle 200. Tại cycle 170, RUL
thật bằng 200 trừ 170, tức 30 cycles. Đơn vị là cycle; dữ liệu không có thông
tin quy đổi ra giờ hoặc ngày, nên em không quy đổi.

Bảng bên phải đối chiếu yêu cầu của đề tài 8 với phần đã làm. Xử lý chuỗi thời
gian nằm ở validation và window theo engine. Về xử lý nhiễu, em dùng thống kê
trên window 30 cycles và Health Index từ PCA; repository chưa dùng một bộ lọc
tín hiệu riêng. Chuẩn hóa và tổng hợp theo window nằm ở slide 10. Trực quan xu
hướng suy giảm ở slide 9 và 11. Dự đoán thời điểm bảo trì và cảnh báo sớm ở
slide 14 và 16.

\slide{4}{Bốn subset C-MAPSS}{50 giây}
\cue{Chỉ vào thẻ FD001, sau đó so sánh ba thẻ còn lại.}
C-MAPSS là dữ liệu mô phỏng động cơ turbofan của NASA. Mỗi subset có file train,
trong đó engine chạy đến failure; file test, trong đó chuỗi bị cắt trước failure;
và file RUL chứa nhãn cho test.

Bốn subset khác nhau theo hai trục: số điều kiện vận hành và số fault mode.
FD001 có một điều kiện và một fault mode là HPC degradation, gồm 100 train, 100
test engines và 33.727 observations. FD002 có sáu điều kiện vận hành. FD003 thêm
fault mode Fan. FD004 có cả sáu điều kiện và hai fault mode.

Mỗi dòng observation có 26 trường: \texttt{engine\_id}, \texttt{cycle}, ba
operational settings và 21 sensor. Trong một subset, khóa là
\texttt{(split, engine\_id, cycle)}. Bốn subset được xử lý riêng vì khóa hiện tại
chưa có \texttt{dataset\_id}.

\slide{5}{Nhãn RUL và ranh giới train/test}{55 giây}
\cue{Đọc hai làn Train và Test từ trái sang phải.}
Ở train, mỗi engine chạy đến failure, nên nhãn suy ra trực tiếp: RUL tại cycle
$t$ bằng cycle cuối trừ $t$. FD001 có 20.631 nhãn train.

Ở test, chuỗi bị cắt trước failure. Không được coi cycle cuối quan sát là
failure. Model tạo đúng một dự đoán tại cycle cuối quan sát của mỗi test engine.
Chỉ sau khi final predictions đã được ghi, pipeline mới ghép file
\texttt{RUL\_FD001.txt} để chấm điểm.

Khung bên phải tóm tắt ba kiểm soát. Nhãn test không dùng trong EDA, đặc trưng
hay chọn model. Nhãn được ghép muộn. Và có dấu vết: \texttt{run\_info.json} ghi
trạng thái \texttt{completed\_after\_final\_predictions}, còn
\texttt{final\_metrics.json} ghi SHA-256 của file nhãn đã dùng. Đây là ranh giới
ngăn data leakage, tức rò rỉ dữ liệu từ test vào quá trình xây dựng mô hình.

\slide{6}{Kiến trúc pipeline}{55 giây}
\cue{Đọc sơ đồ từ trái sang phải, sau đó chỉ hai nhánh bên dưới.}
Bronze là file raw dạng text. Reader kiểm tra đúng 26 trường, kiểu dữ liệu, khóa
duy nhất và cycle liên tục. Train và test cùng schema, nên được thêm cột
\texttt{split} rồi union theo dòng bằng Pandas concat hoặc PySpark
\texttt{unionByName}. Đây là union, không phải join.

Silver lưu \texttt{observations} với 33.727 dòng và \texttt{train\_targets} với
20.631 dòng dưới dạng Parquet. Gold lưu \texttt{engine\_summary} cho 200 engine
dưới dạng Parquet và SQLite. Từ Silver, rolling window 30 tạo 27.927 dòng và
126 cột. Baseline HGB126 dùng 126 cột; predictor HGB169 trong alignment thêm
\texttt{cycle} và 42 history features. Hai biểu diễn được so trên cùng khóa
train và cùng grouped fold.

Hai nhánh bên dưới: EDA, kiểm định thống kê và Health Index chỉ dùng train; nhánh
còn lại xuất metric, hình và dashboard. Khi ghép feature với target, pipeline
dùng đủ ba cột khóa và kiểm tra quan hệ một--một; nếu chỉ join theo
\texttt{engine\_id}, mọi cycle sẽ ghép với mọi cycle và làm nhân dòng.

\slide{7}{Parity Pandas--PySpark và benchmark}{1 phút 5 giây}
\cue{Đọc bảng trái trước, sau đó biểu đồ phải, cuối cùng nêu giới hạn.}
Bảng trái là bằng chứng cho phần Big Data. Cùng một input, pipeline Pandas và
pipeline PySpark cho cùng schema, khóa, số dòng và giá trị ở ba boundary:
Silver observations, Silver train targets và Gold engine summary. Sai khác
tuyệt đối lớn nhất bằng 0.

Biểu đồ phải là benchmark local bổ sung: cùng Parquet input 33.727 và 337.270
dòng, window 30, 126 rolling features. Mỗi cột là trung vị ba lượt đo trong
process riêng, gồm đọc, tổng hợp, rolling và ghi. Startup Spark được báo riêng.
Peak RSS được đo cho cả Python và JVM. Toàn bộ 27.927 dòng rolling feature ở
scale 1 lần đều khớp giữa Pandas và Spark với tolerance $10^{-7}$.

Giới hạn cần nói rõ: Spark chạy local[4]. Output được đưa qua Arrow về driver
rồi ghi bằng PyArrow trong môi trường Windows; write time
bao gồm bước chuyển này. Chưa đo distributed writer hoặc cluster nhiều máy.

\slide{8}{EDA}{30 giây}
\cue{Chỉ vào histogram, sau đó chuyển sang các cột không đổi.}
EDA chỉ dùng train. Tuổi thọ 100 train engines phân tán rộng: ngắn nhất 128,
trung vị 199, trung bình 206,31 và dài nhất 362 cycles, có vài engine dài bất
thường ở đuôi phải.

Bảy cột không đổi trong cả 100 engine: sensor 1, 5, 10, 16, 18, 19 và
\texttt{setting\_3}. Sensor 6 không đổi trong 38 trên 100 engine. Em chưa loại cột
nào: 126 đặc trưng vẫn gồm đủ 21 sensor. Đây là lựa chọn có chủ đích để không dùng
kết quả EDA làm bước chọn biến ngoài training fold.

\slide{9}{Xu hướng suy giảm của sensor}{40 giây}
\cue{Chỉ hai hình trước, sau đó đọc ba bước kiểm định.}
Hai hình là quỹ đạo giá trị gốc của ba train engine có tuổi thọ 178, 199 và 229
cycles. Sensor 4 tăng dần, sensor 7 giảm dần khi engine tiến gần failure. Tín
hiệu có nhiễu rõ theo từng cycle; đây là lý do em tóm tắt theo window ở slide sau.

Để kiểm tra trên toàn bộ train, với mỗi engine em lấy mean 30 cycle cuối trừ mean
30 cycle đầu. Bootstrap theo engine 10.000 lần cho khoảng tin cậy; Wilcoxon
signed-rank cho p-value; Benjamini--Hochberg với $q = 0{,}05$ hiệu chỉnh cho 21
sensor.

Kết quả: 15 trên 21 sensor bác bỏ giả thuyết không sau hiệu chỉnh; 6 sensor hằng
số nên không kiểm định. Phân tích smoothing mới bổ sung dùng trailing mean 30 để
mô tả giảm dao động ngắn hạn; đây là proxy attenuation, không chứng minh đã khử
nhiễu thật hoặc làm RUL tốt hơn. Các kết quả mô tả này không phải feature selection.

\slide{10}{Trailing window và đặc trưng}{45 giây}
\cue{Chỉ dải ô cycle, sau đó sáu thống kê, cuối cùng khung bên phải.}
Với protocol chính, mỗi dòng đặc trưng dùng 30 cycle gần nhất của cùng engine. Ví
dụ tại cycle 170, window gồm cycles 141 đến 170; cycle 171 trở đi không được dùng.

Với mỗi sensor, em tính sáu thống kê: giá trị cuối, mean, sample standard
deviation, min, max và slope tuyến tính. 21 sensor nhân 6 bằng 126 cột. Hai mươi
chín cycle đầu mỗi engine chưa đủ window nên không tạo dòng và không padding. FD001
có 17.731 train windows và 10.196 test windows.

Khung bên phải trả lời yêu cầu về nhiễu và chuẩn hóa. Thống kê trên 30 cycles tóm
tắt dao động từng cycle; repository chưa dùng bộ lọc tín hiệu riêng. ETL chuẩn hóa
schema và kiểu dữ liệu. StandardScaler được fit trong training fold cho Ridge và
trên 90 engine cho Health Index. HGB126/HGB169 dùng feature theo thang gốc; HGB169
thêm cycle và history features, không phải một kiến trúc model khác.

\slide{11}{Health Index và ARIMA}{35 giây}
\cue{Chỉ hình trái để nói về xu hướng suy giảm, sau đó biểu đồ phải.}
Health Index là thành phần chính thứ nhất sau khi chuẩn hóa sensor. Scaler, PCA và
hướng PC1 được fit trên 90 train engines; 10 engine còn lại chỉ được transform.
Hình cho thấy chỉ số giảm dần và giảm nhanh ở cuối vòng đời. Đây là biểu đồ xu
hướng suy giảm mà đề tài yêu cầu.

ARIMA fit 70\% đầu chuỗi và dự báo 30\% cuối. MAE trung bình trên 10 engine là
0,4110, thấp hơn drift 0,7469 và last-value 0,9478, tính trên thang Health Index,
không phải cycle.

Giới hạn: ARIMA chỉ thấp hơn drift ở 7 trên 10 engine, và coverage của khoảng dự
đoán 95\% dao động từ 21,43\% đến 97,03\%. Health Index forecast không phải RUL
forecast; đây là phân tích phụ.

\slide{12}{Evaluation protocol}{45 giây}
\cue{Chỉ lưới fold bên trái, sau đó bốn bước bên phải.}
Em dùng 5-fold GroupKFold với \texttt{engine\_id} làm group. 100 train engines
được chia thành năm nhóm 20 engine; mỗi fold dùng một nhóm làm validation, bốn
nhóm còn lại để fit.

Lý do: hai window liền kề của cùng engine trùng nhau 29 trên 30 cycles. Nếu một
window ở train và window kề nó ở validation, model gần như đã thấy đáp án và metric
sẽ lạc quan.

Bốn bước của protocol gốc là chọn model bằng mean engine-level MAE trên OOF, đặt
sample weight để engine dài không chi phối loss, refit trên train và chỉ ghép nhãn
test sau final prediction. Alignment ngày 07/10 là read-only replay: không fit
model mới, giữ cùng 17.731 khóa, 100 engine, raw truth và grouped fold để so HGB126
với HGB169.

\slide{13}{So sánh HGB126 và HGB169}{55 giây}
\cue{Đọc hai hàng model, sau đó chỉ khung giới hạn.}
HGB126 là baseline lịch sử với 126 rolling features. HGB169 thêm \texttt{cycle} và
42 history features. Trên trajectory panel 17.731 rows, mean engine MAE là
20,2143 với HGB126 và 17,0009 với HGB169; RMSE lần lượt 33,6359 và 29,8068.
Trên endpoint panel ở các fraction 0,5/0,7/0,9, MAE là 15,2577 và 9,8913.

Paired bootstrap theo engine cho thấy chênh lệch có lợi cho HGB169 trên các panel
chung, nhưng khoảng tin cậy có điều kiện trên prediction đã lưu; chưa bao gồm
refit, seed và selection uncertainty. Vì vậy HGB169 là predictor tham chiếu có
điều kiện, còn HGB126 vẫn giữ vai trò baseline audit.

\slide{14}{Metric test lịch sử của HGB126}{1 phút 5 giây}
\cue{Đọc ba thẻ số bên trái, sau đó giải thích hình.}
Đây là metric test lịch sử của HGB126 trên 100 động cơ FD001, không phải một lần
refit mù của HGB169. MAE trung bình theo động cơ là 15,3490 cycles, RMSE là
20,5413 cycles và NASA score là 1.513,2577.

RMSE phạt sai số lớn mạnh hơn MAE. NASA score là score bất đối xứng, phạt dự đoán
trễ nặng hơn và không có đơn vị cycle. Các metric alignment của HGB169 được trình
bày riêng ở slide 13: trajectory MAE 17,0009, RMSE 29,8068; endpoint MAE 9,8913
và RMSE 14,9913. Không trộn hai protocol này khi kết luận.

Đọc hình: trục ngang là RUL thật, trục dọc là RUL dự đoán. Điểm càng gần đường chéo
càng đúng. Vùng RUL thấp bám sát đường chéo hơn vùng RUL cao; đây là vùng quan
trọng cho cảnh báo.

\slide{15}{Phân tích lỗi}{45 giây}
\cue{Chỉ hình residual, sau đó thanh 74/26 và engine 23.}
Hình là residual, tức dự đoán trừ RUL thật, theo RUL thật. Residual dương nghĩa là
model cho rằng engine còn sống lâu hơn thực tế, tương ứng cảnh báo muộn.

Trong view dashboard legacy, có 74 dự đoán cao hơn RUL thật và 26 dự đoán thấp
hơn. Trường hợp sai lớn nhất là engine 23: RUL thật 113, dự đoán 170,4768, sai
57,4768 cycles. Đây là audit case của HGB126, không được gán cho HGB169 nếu chưa
có artifact test tương ứng.

Thông điệp: metric trung bình không thay thế việc xem từng trường hợp sai lớn.

\slide{16}{Cảnh báo sớm và policy}{1 phút 30 giây}
\cue{Đây là sản phẩm đặc thù đề tài yêu cầu. Đọc hình trái trước.}
Policy phải tách khỏi predictor. Với HGB169 và threshold 30 cố định, on-time là
45/100, early-alarm burden 0,0107 và median lead 29 cycles. Threshold 40 cho
96/100 on-time, burden 0,0555 và median lead 38 cycles. Đây là sensitivity policy,
không phải policy triển khai đã được hiệu chỉnh. Calibrated169 historical readout
đạt 94/100 on-time, burden 0,0522 và median lead 37, nhưng dùng inner-selected
policy theo outer fold.

Vì vậy không claim unconditional whole-process improvement: threshold 30 kém
timely hơn với HGB169, còn threshold 40 chỉ nhỉnh hơn về on-time cùng burden tăng
nhẹ. Dashboard legacy HGB126 có bảng cố định 30 với 44 trước, 7 đúng, 49 sau;
đây là số baseline tham chiếu. Mọi train engine đều đi đến failure, nên false-alarm
rate cấp engine không xác định và chưa phải chính sách bảo trì.

\slide{17}{Thử nghiệm cải thiện và giới hạn kết luận}{1 phút 15 giây}
\cue{Chỉ sơ đồ nested CV, bảng kết quả gốc, rồi bảng follow-up phía dưới.}
Sau kết quả baseline, em thử cải thiện vùng RUL thấp và giảm cảnh báo trễ. Em dùng
nested GroupKFold 5 nhân 3: vòng ngoài giữ riêng 20 engine để đánh giá; vòng trong
chia phần còn lại thành ba nhóm để chọn window 10, 20 hoặc 30, một trong 24 cấu hình
HGB và một trong ba cách đặt trọng số. Test không dùng để chọn cấu hình.

Bảng follow-up lịch sử trên cùng 17.731 dòng: MAE khi RUL không quá 30 đổi từ 3,6676 thành
3,6730 cycles, chênh lệch +0,0054 với khoảng tin cậy 95\% [$-0{,}0506$; 0,0639];
cảnh báo trễ hoặc bỏ lỡ tăng từ 49 lên 53. Gate không đạt, nên không promote biến
thể này. Alignment HGB169 là read-only replay riêng, không phải một lần fit mới.

Bảng phía dưới là thí nghiệm bổ sung ngày 29/09: FD001 thử capped RUL và history
features; FD002/FD004 thử settings và chuẩn hóa theo sáu điều kiện vận hành.
Mỗi normalizer chỉ fit trên training engines của từng fold. Cột cuối cho biết
có đạt tiêu chí outer CV đã chốt trước hay không. FD001 giảm MAE RUL thấp từ
3,67 xuống 3,22 nhưng tăng cảnh báo trễ từ 49 lên 55, nên chưa đủ điều kiện thay
baseline. FD002/FD004 là huấn luyện trong từng subset, không phải transfer từ
FD001. Test lịch sử đã được xem nên kết quả follow-up có tính khám phá.

\slide{18}{Chuyển miền FD001--FD004}{55 giây}
\cue{Nói rõ window 10 trước khi đọc biểu đồ.}
Câu hỏi cuối: model FD001 có dùng được cho subset khác không. Em chạy cùng pipeline
với window 10 cycles, vì FD004 có test engine ngắn nhất chỉ 19 cycles. Các số này
không so với MAE 15,3490 dùng window 30.

Cột xanh nhạt là train và test cùng subset: FD002 22,3556, FD003 28,3255 và FD004
27,2218 cycles. Cột cam là giữ model FD001 và áp dụng sang test subset khác: 60,3215,
47,3949 và 58,5731 cycles. MAE tăng từ khoảng 1,7 đến 2,7 lần.

Khác biệt điều kiện vận hành ở FD002 và fault mode ở FD003 đều làm MAE tăng; FD004
có cả hai. Kết luận: cần huấn luyện lại hoặc domain adaptation trước khi dùng model
FD001 cho subset khác. Đây là negative result nhưng trả lời trực tiếp câu hỏi về
khả năng tổng quát hóa.

\slide{19}{Trả lời sáu câu hỏi nghiên cứu}{1 phút 15 giây}
\cue{Đọc sáu thẻ theo hai hàng, sau đó chốt đóng góp Big Data.}
Q1: FD001 đạt kiểm tra schema, khóa và cycle liên tục. Q2: 15 trên 21 sensor
thay đổi có ý nghĩa giữa đầu và cuối vòng đời, và Health Index giảm rõ ở cuối
vòng đời. Q3: PySpark khớp Pandas trên ba bảng Silver/Gold.

Q4: HGB169 là predictor candidate có điều kiện; trajectory MAE 17,0009 so với
20,2143 của HGB126, endpoint MAE 9,8913 so với 15,2577. Bảng test 15,3490 của
HGB126 vẫn được giữ như baseline lịch sử, không trộn với panel alignment.

Q5: với HGB169, threshold 30 có 45\% on-time, còn threshold 40 có 96\% nhưng
burden tăng nhẹ; đây là sensitivity, chưa phải policy triển khai. Q6: model FD001
không chuyển trực tiếp sang FD002--FD004.

Đóng góp Big Data của báo cáo là batch ETL nhiều file, Parquet Silver/Gold, parity
Pandas--PySpark và artifact có hash, \texttt{run\_info} để tái lập.

\slide{20}{Giới hạn, hướng tiếp theo và hỏi đáp}{1 phút}
\cue{Đọc nhanh hai cột, sau đó mời đặt câu hỏi.}
Bốn giới hạn chính: dữ liệu mô phỏng và FD001 chỉ có một điều kiện vận hành, một
fault mode; mỗi test engine chỉ được đánh giá ở một điểm cắt; không có chi phí,
downtime hay intervention records nên chưa tối ưu được lịch bảo trì; PySpark chạy
local; benchmark có I/O nhưng ghi qua driver và chưa đo cluster.

Capped RUL, history features, condition normalization và rolling features bằng
Spark đã được triển khai và kiểm tra. Hướng tiếp theo là xác nhận mù predictor
HGB169, hiệu chỉnh warning policy trên protocol độc lập, cảnh báo theo toàn bộ
quỹ đạo và benchmark distributed writer trên cluster thực.

Em xin cảm ơn thầy/cô đã lắng nghe và xin nhận câu hỏi. Nếu thầy/cô muốn xem trực
tiếp, em có thể mở dashboard theo \filepath{handoff/DEMO_SCRIPT.md}, khoảng ba phút.

\clearpage
\section*{Câu hỏi phản biện thường gặp}
\addcontentsline{toc}{section}{Câu hỏi phản biện thường gặp}

\answer{1. Vì sao chia GroupKFold theo engine?}{Vì các window trong cùng engine có phụ thuộc thời gian và chồng lấn 29/30 cycles. Nếu chia ngẫu nhiên theo row, window rất gần nhau có thể nằm ở cả train và validation, làm metric lạc quan. GroupKFold bảo đảm một engine chỉ thuộc một fold validation.}

\answer{2. Có data leakage không?}{Pipeline không dùng tương lai trong trailing window, giữ mọi window của engine trong cùng fold và chỉ đọc test labels sau khi final predictions đã được ghi. \texttt{run\_info.json} ghi trạng thái join và \texttt{final\_metrics.json} ghi SHA-256 của file nhãn.}

\answer{3. Xử lý nhiễu ở đâu?}{Thống kê mean, min, max, std và slope trên 30 cycles tóm tắt dao động từng cycle; Health Index gom 21 sensor bằng PCA. Repository chưa dùng bộ lọc tín hiệu riêng như moving median hoặc Kalman filter; đây là giới hạn nếu cô hỏi.}

\answer{4. Vì sao không áp dụng model FD001 trực tiếp cho FD002--FD004?}{Kết quả transfer cho thấy MAE tăng từ khoảng 22--28 cycles trong evaluation nội bộ lên 47--60 cycles khi giữ model FD001 cố định. Khác biệt phân bố giữa các subset làm model mất độ chính xác, nên cần retraining hoặc domain adaptation.}

\answer{5. Vì sao window 10 và window 30 không so trực tiếp?}{Window 30 cung cấp lịch sử dài hơn và là protocol chính của FD001. Window 10 được dùng để mọi subset đều có đủ test windows, đặc biệt FD004 có test engine ngắn nhất 19 cycles. Thay đổi window làm thay đổi input và metric.}

\answer{6. ARIMA có dự báo RUL không?}{Không. ARIMA forecast Health Index chuẩn hóa từ PCA. Health Index không phải countdown RUL và chỉ là phân tích phụ.}

\answer{7. Vì sao không dùng LSTM hoặc Transformer?}{Phần thảo luận của báo cáo ghi rằng LSTM/Transformer cần một giả thuyết về biểu diễn thời gian và phải so trên cùng protocol với baseline hiện tại. Báo cáo chưa chạy thí nghiệm đó, nên không đưa ra nhận định về hiệu quả của các model này.}

\answer{8. Kết quả có áp dụng cho động cơ thực tế không?}{Chưa thể kết luận. C-MAPSS là dữ liệu mô phỏng và không có chi phí bảo trì, downtime hoặc intervention records. Cần dữ liệu thực tế và evaluation ngoài miền trước khi triển khai.}

\answer{9. Big Data nằm ở đâu khi dữ liệu gốc nhỏ?}{Đóng góp Big Data nằm ở batch ETL nhiều file, Parquet Silver/Gold, parity Pandas/PySpark, provenance và benchmark local có đọc--rolling--ghi ở scale 1×/10×. Spark chạy local[4] và ghi qua driver; repository không tuyên bố FD001 tự thân là workload production-scale và chưa có benchmark cluster.}

\end{document}
